% =====================================================================
%  Anatomía de un LLM — Fase I — Módulo 2
%  Álgebra lineal I: vectores, matrices y embeddings
%  Compilar: latexmk -r ../latexmkrc M02-algebra-lineal-1.tex
% =====================================================================
\documentclass[11pt,a4paper,twoside,openany]{book}
\usepackage{../estilo-apuntes}

\newcommand{\ApuntesCurso}{Anatomía de un LLM}
\newcommand{\ApuntesModulo}{Módulo 2}

% Notación propia del módulo
\newcommand{\ip}[2]{\langle #1,#2\rangle}
\newcommand{\unos}{\mathbf{1}}
\DeclareMathOperator{\Var}{Var}
\DeclareMathOperator{\diag}{diag}

\hypersetup{pdftitle={Módulo 2. Álgebra lineal I (resumen)},pdfauthor={Anatomía de un LLM}}

\makeatletter
\newlabel{ej:operaciones}{{1.5}{2}{}{ejemplo.5}{}}
\newlabel{ej:operaciones@cref}{{[ejemplo][5][]1.5}{[1][2][]2}}
\newlabel{eje:1-combinacion}{{1.1}{4}{}{ejercicio.1.1}{}}
\newlabel{eje:1-combinacion@cref}{{[ejercicio][1][1]1.1}{[1][4][]4}}
\newlabel{eje:1-recuentos}{{1.2}{4}{}{ejercicio.1.2}{}}
\newlabel{eje:1-recuentos@cref}{{[ejercicio][2][1]1.2}{[1][4][]4}}
\newlabel{eje:1-normas}{{1.3}{4}{}{ejercicio.1.3}{}}
\newlabel{eje:1-normas@cref}{{[ejercicio][3][1]1.3}{[1][4][]4}}
\newlabel{eje:1-cs}{{1.4}{4}{}{ejercicio.1.4}{}}
\newlabel{eje:1-cs@cref}{{[ejercicio][4][1]1.4}{[1][4][]4}}
\newlabel{eje:1-paralelogramo}{{1.5}{4}{}{ejercicio.1.5}{}}
\newlabel{eje:1-paralelogramo@cref}{{[ejercicio][5][1]1.5}{[1][4][]4}}
\newlabel{fig:proyeccion}{{2.1}{6}{Proyección ortogonal de $u$ sobre la recta generada por $a$ (\cref {prop:proyeccion}). El error $u-p$ es perpendicular a $a$}{figure.2.1}{}}
\newlabel{fig:proyeccion@cref}{{[figure][1][2]2.1}{[1][6][]6}}
\newlabel{ej:proyeccion}{{2.6}{6}{}{ejemplo.6}{}}
\newlabel{ej:proyeccion@cref}{{[ejemplo][6][]2.6}{[1][6][]6}}
\newlabel{ej:juguete}{{2.9}{7}{Embeddings de juguete}{ejemplo.9}{}}
\newlabel{ej:juguete@cref}{{[ejemplo][9][]2.9}{[1][7][]7}}
\newlabel{eje:2-escala}{{2.1}{8}{}{ejercicio.2.1}{}}
\newlabel{eje:2-escala@cref}{{[ejercicio][1][2]2.1}{[1][8][]8}}
\newlabel{eje:2-vf}{{2.2}{8}{}{ejercicio.2.2}{}}
\newlabel{eje:2-vf@cref}{{[ejercicio][2][2]2.2}{[1][8][]8}}
\newlabel{eje:2-angulo}{{2.3}{8}{}{ejercicio.2.3}{}}
\newlabel{eje:2-angulo@cref}{{[ejercicio][3][2]2.3}{[1][8][]8}}
\newlabel{eje:2-vecinos}{{2.4}{9}{}{ejercicio.2.4}{}}
\newlabel{eje:2-vecinos@cref}{{[ejercicio][4][2]2.4}{[1][8][]9}}
\newlabel{eje:2-normalizar}{{2.5}{9}{}{ejercicio.2.5}{}}
\newlabel{eje:2-normalizar@cref}{{[ejercicio][5][2]2.5}{[1][9][]9}}
\newlabel{eje:2-tres}{{2.6}{9}{}{ejercicio.2.6}{}}
\newlabel{eje:2-tres@cref}{{[ejercicio][6][2]2.6}{[1][9][]9}}
\newlabel{eje:2-distancia}{{2.7}{9}{}{ejercicio.2.7}{}}
\newlabel{eje:2-distancia@cref}{{[ejercicio][7][2]2.7}{[1][9][]9}}
\newlabel{ej:matvec}{{3.5}{11}{}{ejemplo.5}{}}
\newlabel{ej:matvec@cref}{{[ejemplo][5][]3.5}{[1][11][]11}}
\newlabel{fig:cizalla}{{3.1}{12}{La cizalla $A=\begin {pmatrix}1&1\\0&1\end {pmatrix}$ transforma el cuadrado unidad (discontinuo) en un paralelogramo. La imagen de $e_1$ no cambia y la de $e_2$ es la segunda columna de $A$}{figure.3.1}{}}
\newlabel{fig:cizalla@cref}{{[figure][1][3]3.1}{[1][11][]12}}
\newlabel{eje:3-lineal}{{3.1}{13}{}{ejercicio.3.1}{}}
\newlabel{eje:3-lineal@cref}{{[ejercicio][1][3]3.1}{[1][13][]13}}
\newlabel{eje:3-dos-lecturas}{{3.2}{13}{}{ejercicio.3.2}{}}
\newlabel{eje:3-dos-lecturas@cref}{{[ejercicio][2][3]3.2}{[1][13][]13}}
\newlabel{eje:3-base}{{3.3}{13}{}{ejercicio.3.3}{}}
\newlabel{eje:3-base@cref}{{[ejercicio][3][3]3.3}{[1][13][]13}}
\newlabel{eje:3-rotacion}{{3.4}{13}{}{ejercicio.3.4}{}}
\newlabel{eje:3-rotacion@cref}{{[ejercicio][4][3]3.4}{[1][13][]13}}
\newlabel{eje:3-markov}{{3.5}{13}{}{ejercicio.3.5}{}}
\newlabel{eje:3-markov@cref}{{[ejercicio][5][3]3.5}{[1][13][]13}}
\newlabel{eje:3-rectas}{{3.6}{13}{}{ejercicio.3.6}{}}
\newlabel{eje:3-rectas@cref}{{[ejercicio][6][3]3.6}{[1][13][]13}}
\newlabel{ej:intercambio}{{4.3}{14}{}{ejemplo.3}{}}
\newlabel{ej:intercambio@cref}{{[ejemplo][3][]4.3}{[1][14][]14}}
\newlabel{ej:gram}{{4.11}{16}{}{ejemplo.11}{}}
\newlabel{ej:gram@cref}{{[ejemplo][11][]4.11}{[1][16][]16}}
\newlabel{eje:4-formas}{{4.1}{17}{}{ejercicio.4.1}{}}
\newlabel{eje:4-formas@cref}{{[ejercicio][1][4]4.1}{[1][16][]17}}
\newlabel{eje:4-rs}{{4.2}{17}{}{ejercicio.4.2}{}}
\newlabel{eje:4-rs@cref}{{[ejercicio][2][4]4.2}{[1][17][]17}}
\newlabel{eje:4-ab-ba}{{4.3}{17}{}{ejercicio.4.3}{}}
\newlabel{eje:4-ab-ba@cref}{{[ejercicio][3][4]4.3}{[1][17][]17}}
\newlabel{eje:4-distributiva}{{4.4}{17}{}{ejercicio.4.4}{}}
\newlabel{eje:4-distributiva@cref}{{[ejercicio][4][4]4.4}{[1][17][]17}}
\newlabel{eje:4-coste}{{4.5}{17}{}{ejercicio.4.5}{}}
\newlabel{eje:4-coste@cref}{{[ejercicio][5][4]4.5}{[1][17][]17}}
\newlabel{eje:4-proyeccion}{{4.6}{17}{}{ejercicio.4.6}{}}
\newlabel{eje:4-proyeccion@cref}{{[ejercicio][6][4]4.6}{[1][17][]17}}
\newlabel{eje:4-ortogonal}{{4.7}{17}{}{ejercicio.4.7}{}}
\newlabel{eje:4-ortogonal@cref}{{[ejercicio][7][4]4.7}{[1][17][]17}}
\newlabel{ej:embedding}{{5.4}{18}{}{ejemplo.4}{}}
\newlabel{ej:embedding@cref}{{[ejemplo][4][]5.4}{[1][18][]18}}
\newlabel{ej:atencion}{{5.11}{20}{}{ejemplo.11}{}}
\newlabel{ej:atencion@cref}{{[ejemplo][11][]5.11}{[1][20][]20}}
\newlabel{eje:5-tamano}{{5.1}{21}{}{ejercicio.5.1}{}}
\newlabel{eje:5-tamano@cref}{{[ejercicio][1][5]5.1}{[1][20][]21}}
\newlabel{eje:5-onehot}{{5.2}{21}{}{ejercicio.5.2}{}}
\newlabel{eje:5-onehot@cref}{{[ejercicio][2][5]5.2}{[1][21][]21}}
\newlabel{eje:5-busqueda}{{5.3}{21}{}{ejercicio.5.3}{}}
\newlabel{eje:5-busqueda@cref}{{[ejercicio][3][5]5.3}{[1][21][]21}}
\newlabel{eje:5-analogias}{{5.4}{21}{}{ejercicio.5.4}{}}
\newlabel{eje:5-analogias@cref}{{[ejercicio][4][5]5.4}{[1][21][]21}}
\newlabel{eje:5-secuencia}{{5.5}{21}{}{ejercicio.5.5}{}}
\newlabel{eje:5-secuencia@cref}{{[ejercicio][5][5]5.5}{[1][21][]21}}
\newlabel{eje:5-logits}{{5.6}{21}{}{ejercicio.5.6}{}}
\newlabel{eje:5-logits@cref}{{[ejercicio][6][5]5.6}{[1][21][]21}}
\newlabel{eje:5-atencion}{{5.7}{21}{}{ejercicio.5.7}{}}
\newlabel{eje:5-atencion@cref}{{[ejercicio][7][5]5.7}{[1][21][]21}}
\newlabel{ap:soluciones}{{A}{22}{Soluciones de los ejercicios}{appendix.A}{}}
\newlabel{ap:soluciones@cref}{{[appendix][1][2147483647]A}{[1][22][]22}}
\newlabel{sol:eje:1-combinacion}{{A.1}{22}{Tema 1}{section*.3}{}}
\newlabel{sol:eje:1-combinacion@cref}{{[subappendix][1][2147483647,1]A.1}{[1][22][]22}}
\newlabel{sol:eje:1-recuentos}{{A.1}{22}{Tema 1}{section*.4}{}}
\newlabel{sol:eje:1-recuentos@cref}{{[subappendix][1][2147483647,1]A.1}{[1][22][]22}}
\newlabel{sol:eje:1-normas}{{A.1}{23}{Tema 1}{section*.5}{}}
\newlabel{sol:eje:1-normas@cref}{{[subappendix][1][2147483647,1]A.1}{[1][23][]23}}
\newlabel{sol:eje:1-cs}{{A.1}{23}{Tema 1}{section*.6}{}}
\newlabel{sol:eje:1-cs@cref}{{[subappendix][1][2147483647,1]A.1}{[1][23][]23}}
\newlabel{sol:eje:1-paralelogramo}{{A.1}{23}{Tema 1}{section*.7}{}}
\newlabel{sol:eje:1-paralelogramo@cref}{{[subappendix][1][2147483647,1]A.1}{[1][23][]23}}
\newlabel{sol:eje:2-escala}{{A.2}{24}{Tema 2}{section*.8}{}}
\newlabel{sol:eje:2-escala@cref}{{[subappendix][2][2147483647,1]A.2}{[1][24][]24}}
\newlabel{sol:eje:2-vf}{{A.2}{24}{Tema 2}{section*.9}{}}
\newlabel{sol:eje:2-vf@cref}{{[subappendix][2][2147483647,1]A.2}{[1][24][]24}}
\newlabel{sol:eje:2-angulo}{{A.2}{25}{Tema 2}{section*.10}{}}
\newlabel{sol:eje:2-angulo@cref}{{[subappendix][2][2147483647,1]A.2}{[1][25][]25}}
\newlabel{sol:eje:2-vecinos}{{A.2}{25}{Tema 2}{section*.11}{}}
\newlabel{sol:eje:2-vecinos@cref}{{[subappendix][2][2147483647,1]A.2}{[1][25][]25}}
\newlabel{sol:eje:2-normalizar}{{A.2}{25}{Tema 2}{section*.12}{}}
\newlabel{sol:eje:2-normalizar@cref}{{[subappendix][2][2147483647,1]A.2}{[1][25][]25}}
\newlabel{sol:eje:2-tres}{{A.2}{26}{Tema 2}{section*.13}{}}
\newlabel{sol:eje:2-tres@cref}{{[subappendix][2][2147483647,1]A.2}{[1][26][]26}}
\newlabel{sol:eje:2-distancia}{{A.2}{26}{Tema 2}{section*.14}{}}
\newlabel{sol:eje:2-distancia@cref}{{[subappendix][2][2147483647,1]A.2}{[1][26][]26}}
\newlabel{sol:eje:3-lineal}{{A.3}{27}{Tema 3}{section*.15}{}}
\newlabel{sol:eje:3-lineal@cref}{{[subappendix][3][2147483647,1]A.3}{[1][26][]27}}
\newlabel{sol:eje:3-dos-lecturas}{{A.3}{27}{Tema 3}{section*.16}{}}
\newlabel{sol:eje:3-dos-lecturas@cref}{{[subappendix][3][2147483647,1]A.3}{[1][27][]27}}
\newlabel{sol:eje:3-base}{{A.3}{27}{Tema 3}{section*.17}{}}
\newlabel{sol:eje:3-base@cref}{{[subappendix][3][2147483647,1]A.3}{[1][27][]27}}
\newlabel{sol:eje:3-rotacion}{{A.3}{28}{Tema 3}{section*.18}{}}
\newlabel{sol:eje:3-rotacion@cref}{{[subappendix][3][2147483647,1]A.3}{[1][28][]28}}
\newlabel{sol:eje:3-markov}{{A.3}{28}{Tema 3}{section*.19}{}}
\newlabel{sol:eje:3-markov@cref}{{[subappendix][3][2147483647,1]A.3}{[1][28][]28}}
\newlabel{sol:eje:3-rectas}{{A.3}{28}{Tema 3}{section*.20}{}}
\newlabel{sol:eje:3-rectas@cref}{{[subappendix][3][2147483647,1]A.3}{[1][28][]28}}
\newlabel{sol:eje:4-formas}{{A.4}{29}{Tema 4}{section*.21}{}}
\newlabel{sol:eje:4-formas@cref}{{[subappendix][4][2147483647,1]A.4}{[1][29][]29}}
\newlabel{sol:eje:4-rs}{{A.4}{29}{Tema 4}{section*.22}{}}
\newlabel{sol:eje:4-rs@cref}{{[subappendix][4][2147483647,1]A.4}{[1][29][]29}}
\newlabel{sol:eje:4-ab-ba}{{A.4}{30}{Tema 4}{section*.23}{}}
\newlabel{sol:eje:4-ab-ba@cref}{{[subappendix][4][2147483647,1]A.4}{[1][30][]30}}
\newlabel{sol:eje:4-distributiva}{{A.4}{30}{Tema 4}{section*.24}{}}
\newlabel{sol:eje:4-distributiva@cref}{{[subappendix][4][2147483647,1]A.4}{[1][30][]30}}
\newlabel{sol:eje:4-coste}{{A.4}{31}{Tema 4}{section*.25}{}}
\newlabel{sol:eje:4-coste@cref}{{[subappendix][4][2147483647,1]A.4}{[1][31][]31}}
\newlabel{sol:eje:4-proyeccion}{{A.4}{31}{Tema 4}{section*.26}{}}
\newlabel{sol:eje:4-proyeccion@cref}{{[subappendix][4][2147483647,1]A.4}{[1][31][]31}}
\newlabel{sol:eje:4-ortogonal}{{A.4}{32}{Tema 4}{section*.27}{}}
\newlabel{sol:eje:4-ortogonal@cref}{{[subappendix][4][2147483647,1]A.4}{[1][31][]32}}
\newlabel{sol:eje:5-tamano}{{A.5}{32}{Tema 5}{section*.28}{}}
\newlabel{sol:eje:5-tamano@cref}{{[subappendix][5][2147483647,1]A.5}{[1][32][]32}}
\newlabel{sol:eje:5-onehot}{{A.5}{33}{Tema 5}{section*.29}{}}
\newlabel{sol:eje:5-onehot@cref}{{[subappendix][5][2147483647,1]A.5}{[1][33][]33}}
\newlabel{sol:eje:5-busqueda}{{A.5}{33}{Tema 5}{section*.30}{}}
\newlabel{sol:eje:5-busqueda@cref}{{[subappendix][5][2147483647,1]A.5}{[1][33][]33}}
\newlabel{sol:eje:5-analogias}{{A.5}{33}{Tema 5}{section*.31}{}}
\newlabel{sol:eje:5-analogias@cref}{{[subappendix][5][2147483647,1]A.5}{[1][33][]33}}
\newlabel{sol:eje:5-secuencia}{{A.5}{34}{Tema 5}{section*.32}{}}
\newlabel{sol:eje:5-secuencia@cref}{{[subappendix][5][2147483647,1]A.5}{[1][34][]34}}
\newlabel{sol:eje:5-logits}{{A.5}{34}{Tema 5}{section*.33}{}}
\newlabel{sol:eje:5-logits@cref}{{[subappendix][5][2147483647,1]A.5}{[1][34][]34}}
\newlabel{sol:eje:5-atencion}{{A.5}{35}{Tema 5}{section*.34}{}}
\newlabel{sol:eje:5-atencion@cref}{{[subappendix][5][2147483647,1]A.5}{[1][35][]35}}
\makeatother

\begin{document}
\frontmatter
\begin{titlepage}
\centering
\vspace*{2.5cm}
{\large\scshape Anatomía de un LLM}\par
\vspace{0.4cm}
{\normalsize Fundamentos matemáticos de los modelos de lenguaje}\par
\vspace{2.2cm}
\rule{\textwidth}{0.8pt}\par\vspace{0.5cm}
{\small\scshape Fase I \textperiodcentered\ Fundamentos matemáticos}\par\vspace{0.3cm}
{\Huge\bfseries Módulo 2\par}\vspace{0.3cm}
{\LARGE Álgebra lineal I\par}
\vspace{0.2cm}
{\large Vectores, matrices y embeddings\par}
\vspace{0.5cm}\rule{\textwidth}{0.4pt}\par
\vspace{1.6cm}
{\large Resumen del módulo}\par
\vfill
{\small Edición de octubre de 2026}\par
\end{titlepage}

\chapter*{Sobre este resumen}
\addcontentsline{toc}{chapter}{Sobre este resumen}
\markboth{Sobre este resumen}{}

Este documento reúne lo indispensable de los apuntes del Módulo 2: el texto expositivo, todas las definiciones y todos los resultados (teoremas, proposiciones, lemas y corolarios), las tablas y las observaciones de contenido. Se han suprimido las demostraciones, los ejemplos, las lecturas, las figuras, los ejercicios y sus soluciones. La numeración de temas, secciones, enunciados y ecuaciones es la misma que en los apuntes, de modo que cualquier referencia puede buscarse allí; las que remiten a elementos suprimidos (un ejemplo, una figura, una ecuación de una demostración) conservan su número en los apuntes.

\section*{Notación}
\begin{center}
\begin{tabular}{@{}ll@{}}
\toprule
Símbolo & Significado \\
\midrule
$u=(u_1,\dots,u_d)\in\R^{d}$ & vector columna de $d$ componentes \\
$e_1,\dots,e_d$ & vectores de la base canónica (o vectores \emph{one-hot}) \\
$u\cdot v=u^{\top}v$ & producto escalar \\
$\norm{u}$, $\norm{u}_1$, $\norm{u}_\infty$ & normas euclídea, $\ell_1$ y del máximo \\
$\hat u=u/\norm{u}$ & vector normalizado ($u\neq 0$) \\
$\cos(u,v)$ & similitud coseno \\
$A\in\R^{m\times n}$, $A_{ij}$ & matriz de $m$ filas y $n$ columnas y su entrada $(i,j)$ \\
$A^{\top}$, $I_n$ & traspuesta e identidad de orden $n$ \\
$\V$, $K=\abs{\V}$, $E\in\R^{K\times d}$ & vocabulario, su tamaño y matriz de embeddings \\
$q$, $k$, $d_k$ & consulta, clave y su dimensión (atención) \\
\bottomrule
\end{tabular}
\end{center}

Los vectores se escriben con letras minúsculas y se consideran columnas; cuando se escriben en línea, $(u_1,\dots,u_d)$ denota la columna correspondiente. Los índices de los tokens empiezan en~0, como en el módulo~1; los de las componentes de un vector, en~1.


\tableofcontents

\mainmatter
% =====================================================================
\chapter{Vectores, producto escalar y norma}
\label{tema:1}
% =====================================================================

Este tema introduce el espacio $\R^{d}$, sus dos operaciones básicas y las dos herramientas métricas sobre las que descansa todo el módulo: el producto escalar y la norma. La presentación sigue las secciones~1.1 y~1.2 de Strang~\cite{strang} y la clase~1 del curso 18.06SC~\cite{mit1806sc}; las definiciones abstractas de espacio vectorial, norma y producto interior proceden de~\cite[§2.4, §3.1 y §3.2]{mml}, y el punto de vista computacional, de~\cite[cap.~1 y §3.1]{vmls}.

\section{El espacio \texorpdfstring{$\R^{d}$}{ℝᵈ}}

\begin{definicion}[Vector]
Sea $d\ge 1$ un entero. Un \emph{vector} de $\R^{d}$ es una lista ordenada $u=(u_1,\dots,u_d)$ de $d$ números reales, llamados \emph{componentes} de $u$. Dos vectores son iguales si lo son todas sus componentes. El \emph{vector nulo} es $0=(0,\dots,0)$.
\end{definicion}

Un vector de $\R^{2}$ o de $\R^{3}$ admite dos lecturas geométricas equivalentes: un punto del plano o del espacio, o una flecha que parte del origen y termina en ese punto. Para $d>3$ no hay dibujo posible, pero todas las definiciones de este módulo se formulan componente a componente y valen igual en cualquier dimensión.

\begin{definicion}[Suma y producto por un escalar]
Para $u,v\in\R^{d}$ y $c\in\R$ se definen
\[
u+v=(u_1+v_1,\dots,u_d+v_d),\qquad c\,u=(c\,u_1,\dots,c\,u_d).
\]
\end{definicion}

Geométricamente, $u+v$ se obtiene colocando la flecha $v$ a continuación de la flecha $u$ (regla del paralelogramo), y $c\,u$ estira la flecha si $c>1$, la encoge si $0<c<1$ y la invierte si $c<0$. La diferencia $u-v=u+(-1)v$ es la flecha que va del extremo de $v$ al extremo de $u$.

\begin{proposicion}\label{prop:ev}
Con la suma y el producto por un escalar, $\R^{d}$ satisface, para todos $u,v,w\in\R^{d}$ y $a,b\in\R$:
\begin{enumerate}
  \item $u+v=v+u$ y $(u+v)+w=u+(v+w)$;
  \item $u+0=u$ y $u+(-u)=0$;
  \item $a(u+v)=au+av$, $(a+b)u=au+bu$ y $(ab)u=a(bu)$;
  \item $1\,u=u$.
\end{enumerate}
\end{proposicion}

Un conjunto dotado de una suma y de un producto por escalares reales que cumplen las propiedades de la \cref{prop:ev} se llama \emph{espacio vectorial real}~\cite[§2.4]{mml}. Además de $\R^{d}$, son espacios vectoriales el conjunto $\R^{m\times n}$ de matrices (\cref{tema:3}) y el de polinomios de grado a lo sumo $n$. Este módulo trabaja casi siempre en $\R^{d}$.

\begin{definicion}[Combinación lineal]
Sean $v_1,\dots,v_k\in\R^{d}$ y $c_1,\dots,c_k\in\R$. El vector $c_1v_1+\dots+c_kv_k$ es una \emph{combinación lineal} de $v_1,\dots,v_k$ con coeficientes $c_1,\dots,c_k$.
\end{definicion}

\addtocounter{definicion}{1}%

Las combinaciones lineales son la operación central del álgebra lineal~\cite[§1.1]{strang}: el producto de una matriz por un vector (\cref{tema:3}) es una combinación lineal de las columnas de la matriz, y la cuestión de qué vectores se obtienen como combinación de otros dados conduce a los conceptos de independencia, base y rango del módulo~3.

\begin{observacion}[Vectores en los modelos de lenguaje]
Un modelo de lenguaje asigna a cada token un vector de $\R^{d}$, su \emph{embedding} (\cref{tema:5}). En el modelo más pequeño de GPT-2 es $d=768$~\cite{radford}, y en GPT-3, con 175\,000 millones de parámetros, $d=12\,288$~\cite[tabla~2.1]{brown}. Anthropic, Google y xAI no publican este valor para los modelos actuales de Claude, Gemini y Grok. Otro ejemplo clásico es el \emph{vector de recuentos} de un documento: fijado un diccionario de $n$ palabras, la componente $i$ es el número de veces que aparece la palabra $i$~\cite[§1.1]{vmls}.
\end{observacion}

\section{El producto escalar}

\begin{definicion}[Producto escalar]\label{def:escalar}
El \emph{producto escalar} (o producto interior euclídeo) de $u,v\in\R^{d}$ es el número real
\[
u\cdot v=\sum_{i=1}^{d}u_iv_i.
\]
Con la notación matricial del \cref{tema:3}, $u\cdot v=u^{\top}v$.
\end{definicion}

\addtocounter{definicion}{1}%

\begin{proposicion}[Propiedades del producto escalar]\label{prop:escalar}
Para todos $u,v,w\in\R^{d}$ y $a,b\in\R$:
\begin{enumerate}
  \item (simetría) $u\cdot v=v\cdot u$;
  \item (linealidad en cada argumento) $(au+bv)\cdot w=a\,(u\cdot w)+b\,(v\cdot w)$;
  \item (positividad) $u\cdot u\ge 0$, con igualdad si y solo si $u=0$.
\end{enumerate}
\end{proposicion}

Una aplicación $\ip{\cdot}{\cdot}\colon V\times V\to\R$ sobre un espacio vectorial que cumple (i)--(iii) se llama \emph{producto interior}~\cite[§3.2]{mml}. El producto escalar es el producto interior estándar de $\R^{d}$; todos los resultados de este tema, salvo los que usan componentes, valen para un producto interior cualquiera.

\begin{observacion}[Coste]
Calcular $u\cdot v$ en $\R^{d}$ requiere $d$ multiplicaciones y $d-1$ sumas, unas $2d$ operaciones en coma flotante~\cite[§1.5]{vmls}. Un LLM calcula miles de millones de productos escalares por cada token generado, y su organización en productos de matrices (\cref{tema:4}) es lo que permite ejecutarlos eficientemente en una GPU.
\end{observacion}

\section{Normas}

\begin{definicion}[Norma euclídea]\label{def:norma}
La \emph{norma} (o longitud) euclídea de $u\in\R^{d}$ es
\[
\norm{u}=\sqrt{u\cdot u}=\sqrt{u_1^{2}+\dots+u_d^{2}}.
\]
Un vector $u$ es \emph{unitario} si $\norm{u}=1$. Si $u\neq 0$, el vector $\hat u=u/\norm{u}$ es unitario y tiene la misma dirección y sentido que $u$; se dice que $\hat u$ es el \emph{normalizado} de $u$. La \emph{distancia} entre $u$ y $v$ es $\norm{u-v}$.
\end{definicion}

En $\R^{2}$ y $\R^{3}$ la norma es la longitud de la flecha, por el teorema de Pitágoras. Por ejemplo, $\norm{(3,4)}=5$, $\norm{(1,2,2)}=3$ y el normalizado de $(3,4)$ es $(0{,}6;\,0{,}8)$.

\begin{definicion}[Norma]\label{def:norma-general}
Una \emph{norma} sobre un espacio vectorial $V$ es una aplicación $\norm{\cdot}\colon V\to\R$ que cumple, para todos $u,v\in V$ y $c\in\R$~\cite[§3.1]{mml}:
\begin{enumerate}
  \item (definida positiva) $\norm{u}\ge 0$, con igualdad si y solo si $u=0$;
  \item (homogeneidad) $\norm{c\,u}=\abs{c}\,\norm{u}$;
  \item (desigualdad triangular) $\norm{u+v}\le\norm{u}+\norm{v}$.
\end{enumerate}
\end{definicion}

\addtocounter{definicion}{1}%

Las propiedades (i) y (ii) de la \cref{def:norma-general} son inmediatas para la norma euclídea. La desigualdad triangular, en cambio, requiere un resultado previo que es probablemente la desigualdad más importante del álgebra lineal.

\begin{teorema}[Desigualdad de Cauchy--Schwarz]\label{teo:cs}
Para todos $u,v\in\R^{d}$,
\[
\abs{u\cdot v}\le\norm{u}\,\norm{v},
\]
y la igualdad se da si y solo si uno de los dos vectores es múltiplo del otro.
\end{teorema}

La elección de $t^{*}$ no es arbitraria: $t^{*}v$ es la proyección de $u$ sobre la recta de $v$ (\cref{prop:proyeccion}), y la demostración dice que el cuadrado de la distancia de $u$ a esa recta es no negativo. El ejercicio~\ref{eje:1-cs} propone una demostración alternativa.

\begin{corolario}[Desigualdad triangular]\label{cor:triangular}
Para todos $u,v\in\R^{d}$, $\norm{u+v}\le\norm{u}+\norm{v}$. En consecuencia, la norma euclídea es una norma en el sentido de la \cref{def:norma-general}.
\end{corolario}

\begin{observacion}
La desigualdad triangular expresa que el camino directo es el más corto: la longitud del lado $u+v$ de un triángulo no supera la suma de las de los otros dos. Aplicada a $u-w=(u-v)+(v-w)$, da $\norm{u-w}\le\norm{u-v}+\norm{v-w}$, que es la desigualdad triangular para distancias.
\end{observacion}


% =====================================================================
\chapter{Ángulos, ortogonalidad y similitud coseno}
\label{tema:2}
% =====================================================================

La desigualdad de Cauchy--Schwarz permite definir el ángulo entre dos vectores en cualquier dimensión. Este tema desarrolla esa geometría: ortogonalidad, proyección sobre una recta y similitud coseno, que es la medida con la que se comparan los embeddings. Se sigue la sección~1.2 de Strang~\cite{strang} y las secciones~3.3, 3.4 y~3.8 de~\cite{mml}; el comportamiento en alta dimensión procede del capítulo~2 de Blum, Hopcroft y Kannan~\cite{bhk}.

\section{El ángulo entre dos vectores}

\begin{definicion}[Ángulo]\label{def:angulo}
Sean $u,v\in\R^{d}$ no nulos. El \emph{ángulo} entre $u$ y $v$ es el único $\theta\in[0,\pi]$ tal que
\[
\cos\theta=\frac{u\cdot v}{\norm{u}\,\norm{v}}.
\]
\end{definicion}

La definición es correcta porque, por el \cref{teo:cs}, el cociente está en $[-1,1]$, y la función coseno es una biyección de $[0,\pi]$ en $[-1,1]$. La proposición siguiente muestra que, en el plano, coincide con el ángulo de la geometría elemental.

\begin{proposicion}\label{prop:ley-coseno}
Sean $u,v\in\R^{2}$ no nulos y sea $\varphi$ el ángulo geométrico que forman las flechas $u$ y $v$. Entonces $u\cdot v=\norm{u}\norm{v}\cos\varphi$.
\end{proposicion}

El signo del producto escalar indica, por tanto, si dos vectores apuntan en direcciones parecidas u opuestas (\cref{tab:signo}). El ejercicio~\ref{eje:2-tres} explora hasta dónde pueden «llevarse mal» varios vectores del plano.

\begin{table}[h]
\centering\small
\caption{Signo del producto escalar y ángulo entre dos vectores no nulos.}\label{tab:signo}
\begin{tabular}{@{}lll@{}}
\toprule
Producto escalar & Ángulo & Interpretación \\
\midrule
$u\cdot v>0$ & $0\le\theta<\pi/2$ & direcciones parecidas \\
$u\cdot v=0$ & $\theta=\pi/2$ & direcciones perpendiculares \\
$u\cdot v<0$ & $\pi/2<\theta\le\pi$ & direcciones opuestas \\
\bottomrule
\end{tabular}
\end{table}

\begin{definicion}[Ortogonalidad]
Dos vectores $u,v\in\R^{d}$ son \emph{ortogonales}, y se escribe $u\perp v$, si $u\cdot v=0$.
\end{definicion}

El vector nulo es ortogonal a todos, y es el único vector ortogonal a sí mismo, por la positividad del producto escalar. Los vectores de la base canónica son ortogonales dos a dos: $e_i\cdot e_j=0$ si $i\neq j$.

\begin{teorema}[Pitágoras]\label{teo:pitagoras}
Si $u\perp v$, entonces $\norm{u+v}^{2}=\norm{u}^{2}+\norm{v}^{2}$.
\end{teorema}

\section{Proyección sobre una recta}

\begin{proposicion}[Proyección ortogonal sobre una recta]\label{prop:proyeccion}
Sean $u,a\in\R^{d}$ con $a\neq 0$. Entre todos los vectores de la recta $\{ta: t\in\R\}$, el más próximo a $u$ es
\[
p=\frac{u\cdot a}{a\cdot a}\,a,
\]
y es el único punto de la recta tal que $u-p\perp a$. El número $(u\cdot a)/\norm{a}$ se llama \emph{proyección escalar} de $u$ sobre $a$; si $\norm{a}=1$, coincide con $u\cdot a$.
\end{proposicion}

\addtocounter{figure}{1}%

\addtocounter{definicion}{1}%

La proyección sobre una recta es el caso más sencillo de la proyección sobre un subespacio, que conduce a los mínimos cuadrados~\cite[§4.2]{strang}; en el \cref{tema:4} se escribe como producto de matrices (ejercicio~\ref{eje:4-proyeccion}).

\section{Similitud coseno}

\begin{definicion}[Similitud coseno]\label{def:coseno}
La \emph{similitud coseno} de dos vectores no nulos $u,v\in\R^{d}$ es
\[
\cos(u,v)=\frac{u\cdot v}{\norm{u}\,\norm{v}}=\hat u\cdot\hat v\in[-1,1].
\]
\end{definicion}

\begin{proposicion}\label{prop:coseno}
Sean $u,v\in\R^{d}$ no nulos.
\begin{enumerate}
  \item $\cos(au,bv)=\cos(u,v)$ para todos $a,b>0$, y $\cos(u,-v)=-\cos(u,v)$.
  \item $\cos(u,v)=1$ si y solo si $u=cv$ con $c>0$; $\cos(u,v)=-1$ si y solo si $u=cv$ con $c<0$.
\end{enumerate}
\end{proposicion}

La similitud coseno mide, por tanto, solo la dirección: $\cos((1,2),(2,4))=1$ aunque los vectores tengan longitudes distintas.

\addtocounter{definicion}{1}%

\begin{proposicion}[Coseno y distancia]\label{prop:coseno-distancia}
Si $\norm{u}=\norm{v}=1$, entonces $\norm{u-v}^{2}=2-2\cos(u,v)$. En consecuencia, entre vectores normalizados, ordenar los candidatos por distancia creciente a una consulta equivale a ordenarlos por similitud coseno decreciente.
\end{proposicion}

La demostración se propone en el ejercicio~\ref{eje:2-distancia}. Sin normalizar, ambas medidas pueden discrepar: un vector muy largo en la misma dirección que la consulta tiene coseno~1 pero puede estar lejos (ejercicio~\ref{eje:2-vecinos}). Los buscadores vectoriales de los sistemas RAG (módulo~22) usan el coseno o el producto escalar sobre vectores normalizados.

\section{Geometría en alta dimensión}

En $\R^{2}$, dos direcciones elegidas al azar forman ángulos de cualquier tamaño. En dimensión alta el comportamiento es muy distinto: dos direcciones aleatorias son casi ortogonales~\cite[§2.4.2]{bhk}. El resultado siguiente lo cuantifica.

\begin{proposicion}\label{prop:alta-dim}
Sea $U$ un vector aleatorio de $\R^{d}$ con $\norm{U}=1$ cuya distribución no cambia al permutar sus componentes ni al cambiar el signo de cualquiera de ellas (por ejemplo, $U=G/\norm{G}$ con $G_1,\dots,G_d$ normales estándar independientes). Para todo vector unitario fijo $v$,
\[
\E[U\cdot v]=0,\qquad \E\bigl[(U\cdot v)^{2}\bigr]=\frac1d .
\]
En particular, la similitud coseno entre $U$ y $v$ tiene media~0 y desviación típica $1/\sqrt d$.
\end{proposicion}

Si $v$ también es aleatorio e independiente de $U$, el resultado se aplica condicionando al valor de $v$. Para $d=768$, la desviación típica del coseno entre dos direcciones aleatorias es $1/\sqrt{768}\approx 0{,}0361$; para $d=12\,288$, $\approx 0{,}0090$. Una simulación con 2000 pares de vectores gaussianos da $0{,}0363$ y $0{,}0091$, respectivamente. Por eso un coseno de $0{,}3$ entre dos embeddings de dimensión 768, que equivale a más de ocho desviaciones típicas, es una señal fuerte de parecido.

\begin{observacion}
Este fenómeno permite que un espacio de dimensión $d$ contenga muchas más de $d$ direcciones casi ortogonales entre sí. La hipótesis de \emph{superposición} de la interpretabilidad (módulo~23) se apoya en esta propiedad: un modelo puede representar más conceptos que dimensiones tiene su espacio, asignando a cada uno una dirección casi ortogonal a las demás.
\end{observacion}


% =====================================================================
\chapter{Matrices y aplicaciones lineales}
\label{tema:3}
% =====================================================================

Una matriz es una tabla de números, pero su papel en el álgebra lineal es el de una función: transforma vectores en vectores respetando las combinaciones lineales. Este tema establece esa correspondencia entre matrices y aplicaciones lineales. Se sigue la sección~2.1 de Strang~\cite{strang} y la sesión «The Geometry of Linear Equations» de 18.06SC~\cite{mit1806sc}; las definiciones proceden de~\cite[§2.2 y §2.7]{mml} y de~\cite[cap.~6 y §8.1]{vmls}.

\section{Matrices}

\begin{definicion}[Matriz]
Sean $m,n\ge 1$. Una \emph{matriz} real de tamaño $m\times n$ es una tabla $A=(A_{ij})$ de $mn$ números reales dispuestos en $m$ filas y $n$ columnas; $A_{ij}$ es la entrada de la fila $i$ y la columna $j$. El conjunto de estas matrices se denota $\R^{m\times n}$. Se escribe $a_j\in\R^{m}$ para la columna $j$ y $A=[\,a_1\ \cdots\ a_n\,]$.
\end{definicion}

Un vector de $\R^{d}$ se identifica con una matriz $d\times 1$ (vector columna). Las matrices del mismo tamaño se suman y se multiplican por escalares entrada a entrada, y con esas operaciones $\R^{m\times n}$ es un espacio vectorial.

\begin{definicion}[Traspuesta e identidad]
La \emph{traspuesta} de $A\in\R^{m\times n}$ es la matriz $A^{\top}\in\R^{n\times m}$ con $(A^{\top})_{ij}=A_{ji}$: las filas de $A$ son las columnas de $A^{\top}$. Una matriz cuadrada es \emph{simétrica} si $A^{\top}=A$. La \emph{identidad} $I_n\in\R^{n\times n}$ tiene unos en la diagonal y ceros fuera de ella; sus columnas son los vectores $e_1,\dots,e_n$ de la base canónica.
\end{definicion}

Si $u\in\R^{d}$ es una columna, $u^{\top}$ es una matriz $1\times d$ (vector fila). Esta notación permite escribir el producto escalar como $u\cdot v=u^{\top}v$.

\section{El producto de una matriz por un vector}

\begin{definicion}[Producto matriz-vector]\label{def:matvec}
Sean $A\in\R^{m\times n}$ y $x\in\R^{n}$. El producto $Ax\in\R^{m}$ es el vector de componentes
\[
(Ax)_i=\sum_{j=1}^{n}A_{ij}x_j,\qquad i=1,\dots,m.
\]
El producto solo está definido si el número de columnas de $A$ coincide con el número de componentes de $x$.
\end{definicion}

\begin{proposicion}[Dos lecturas del producto]\label{prop:filas-columnas}
Sea $A\in\R^{m\times n}$ con filas $r_1^{\top},\dots,r_m^{\top}$ y columnas $a_1,\dots,a_n$. Para todo $x\in\R^{n}$:
\begin{enumerate}
  \item (por filas) la componente $i$ de $Ax$ es el producto escalar $r_i\cdot x$;
  \item (por columnas) $Ax=x_1a_1+x_2a_2+\dots+x_na_n$ es una combinación lineal de las columnas de $A$.
\end{enumerate}
\end{proposicion}

\addtocounter{definicion}{1}%

Las dos lecturas corresponden a las dos imágenes de un sistema de ecuaciones lineales $Ax=b$ que Strang llama imagen por filas e imagen por columnas~\cite[§2.1]{strang}: por filas, cada ecuación es un hiperplano y la solución es su intersección; por columnas, se buscan los coeficientes con los que las columnas de $A$ se combinan para dar $b$.

\section{Aplicaciones lineales}

\begin{definicion}[Aplicación lineal]\label{def:lineal}
Una aplicación $T\colon\R^{n}\to\R^{m}$ es \emph{lineal} si
\[
T(au+bv)=a\,T(u)+b\,T(v)\qquad\text{para todos } u,v\in\R^{n},\ a,b\in\R.
\]
\end{definicion}

Tomando $a=b=0$ se obtiene $T(0)=0$: una aplicación lineal fija el origen. Por inducción, $T$ transforma cualquier combinación lineal en la combinación lineal de las imágenes con los mismos coeficientes.

\begin{proposicion}\label{prop:matriz-lineal}
Para toda $A\in\R^{m\times n}$, la aplicación $x\mapsto Ax$ es lineal.
\end{proposicion}

El resultado recíproco es el que justifica la frase «una matriz es una aplicación lineal».

\begin{teorema}[Matriz de una aplicación lineal]\label{teo:representacion}
Sea $T\colon\R^{n}\to\R^{m}$ lineal y sea $A\in\R^{m\times n}$ la matriz cuyas columnas son $T(e_1),\dots,T(e_n)$. Entonces $T(x)=Ax$ para todo $x\in\R^{n}$, y $A$ es la única matriz con esta propiedad.
\end{teorema}

El \cref{teo:representacion} da una receta práctica: para conocer una aplicación lineal basta saber adónde envía los vectores de la base canónica, y esas imágenes son las columnas de su matriz. El \cref{tab:transformaciones} recoge las transformaciones del plano más habituales, y la \cref{fig:cizalla} muestra el efecto de una cizalla sobre el cuadrado unidad.

\begin{table}[h]
\centering\small
\caption{Transformaciones lineales del plano. Las columnas de cada matriz son las imágenes de $e_1$ y $e_2$.}\label{tab:transformaciones}
\begin{tabular}{@{}lccc@{}}
\toprule
Transformación & Matriz & $e_1\mapsto$ & $e_2\mapsto$ \\
\midrule
Rotación de ángulo $\theta$ & $R_\theta=\begin{pmatrix}\cos\theta&-\sin\theta\\ \sin\theta&\cos\theta\end{pmatrix}$ & $(\cos\theta,\sin\theta)$ & $(-\sin\theta,\cos\theta)$ \\[10pt]
Rotación de $90^\circ$ & $\begin{pmatrix}0&-1\\1&0\end{pmatrix}$ & $(0,1)$ & $(-1,0)$ \\[8pt]
Escalado & $\begin{pmatrix}2&0\\0&0{,}5\end{pmatrix}$ & $(2,0)$ & $(0;\,0{,}5)$ \\[8pt]
Cizalla & $\begin{pmatrix}1&1\\0&1\end{pmatrix}$ & $(1,0)$ & $(1,1)$ \\[8pt]
Simetría respecto de $y=x$ & $\begin{pmatrix}0&1\\1&0\end{pmatrix}$ & $(0,1)$ & $(1,0)$ \\[8pt]
Proyección sobre el eje $x$ & $\begin{pmatrix}1&0\\0&0\end{pmatrix}$ & $(1,0)$ & $(0,0)$ \\
\bottomrule
\end{tabular}
\end{table}

\addtocounter{figure}{1}%

\begin{proposicion}[Las rotaciones conservan la norma]\label{prop:rotacion}
Para todo $\theta\in\R$ y todo $x\in\R^{2}$, $\norm{R_\theta x}=\norm{x}$.
\end{proposicion}

Las rotaciones conservan también los productos escalares y, por tanto, los ángulos (ejercicio~\ref{eje:4-ortogonal}). Los escalados no: con la matriz de escalado del \cref{tab:transformaciones}, $e_1+e_2$, de norma $\sqrt2$, pasa a $(2;\,0{,}5)$, de norma $\sqrt{4{,}25}$.

\begin{observacion}[Matrices en los modelos de lenguaje]\label{obs:qkv}
Casi todo el cálculo de un Transformer consiste en aplicar matrices aprendidas a vectores. En la atención (módulo~13), cada embedding $x\in\R^{d}$ se transforma en una \emph{consulta} $q=W_Qx$, una \emph{clave} $k=W_Kx$ y un \emph{valor} $v=W_Vx$, con $W_Q,W_K,W_V$ matrices $d_k\times d$. Estas matrices seleccionan qué rasgos del embedding intervienen en cada comparación.
\end{observacion}

\begin{observacion}[Matrices de transición]
La matriz de transición $P$ de un modelo de bigramas (módulo~1, tema~2) es estocástica por filas. Si $\pi_t\in\R^{K}$ es la distribución del token en la posición $t$, la de la posición siguiente es $\pi_{t+1}=P^{\top}\pi_t$, cuya componente $j$ es $\sum_iP_{ij}(\pi_t)_i$. La matriz $P^{\top}$ tiene columnas que suman~1; es la convención que usa Strang para las \emph{matrices de Markov} (ejercicio~\ref{eje:3-markov}).
\end{observacion}


% =====================================================================
\chapter{El producto de matrices}
\label{tema:4}
% =====================================================================

El producto de matrices se define de modo que represente la composición de aplicaciones lineales. De esa definición se siguen sus propiedades (asociativa, distributivas, comportamiento con la traspuesta) y su principal diferencia con el producto de números: no es conmutativo. El tema sigue la sección~2.4 de Strang~\cite{strang} y la sesión «Multiplication and Inverse Matrices» de 18.06SC~\cite{mit1806sc}; la composición y el coste computacional se tratan como en~\cite[§10.1 y §10.2]{vmls}.

\section{Definición y composición}

\begin{definicion}[Producto de matrices]\label{def:producto}
Sean $A\in\R^{m\times n}$ y $B\in\R^{n\times p}$. El producto $AB\in\R^{m\times p}$ es la matriz de entradas
\[
(AB)_{ij}=\sum_{k=1}^{n}A_{ik}B_{kj},\qquad 1\le i\le m,\ 1\le j\le p.
\]
La entrada $(i,j)$ es el producto escalar de la fila $i$ de $A$ por la columna $j$ de $B$. El producto solo está definido si el número de columnas de $A$ coincide con el número de filas de $B$: $(m\times n)(n\times p)=m\times p$.
\end{definicion}

\begin{teorema}[El producto es la composición]\label{teo:composicion}
Sean $A\in\R^{m\times n}$ y $B\in\R^{n\times p}$. Para todo $x\in\R^{p}$, $(AB)x=A(Bx)$. Es decir, la matriz de la aplicación «aplicar primero $B$ y después $A$» es $AB$.
\end{teorema}

El orden de escritura es el de la composición de funciones, $(A\circ B)(x)=A(B(x))$: en $AB$ actúa primero la matriz de la derecha.

\addtocounter{definicion}{1}%

\section{Cuatro formas de calcular un producto}

La \cref{def:producto} calcula el producto entrada a entrada. Strang subraya otras tres formas equivalentes~\cite[§2.4]{strang}, cada una útil en contextos distintos.

\begin{proposicion}\label{prop:cuatro}
Sean $A\in\R^{m\times n}$ con columnas $a_1,\dots,a_n$ y filas $r_1^{\top},\dots,r_m^{\top}$, y $B\in\R^{n\times p}$ con columnas $b_1,\dots,b_p$ y filas $s_1^{\top},\dots,s_n^{\top}$. Entonces:
\begin{enumerate}
  \item (por columnas) la columna $j$ de $AB$ es $Ab_j$;
  \item (por filas) la fila $i$ de $AB$ es $r_i^{\top}B$;
  \item (columnas por filas) $AB=\sum_{k=1}^{n}a_k\,s_k^{\top}$, suma de $n$ matrices $m\times p$.
\end{enumerate}
\end{proposicion}

Una matriz de la forma $a\,s^{\top}$, con $a\in\R^{m}$ y $s\in\R^{p}$, se llama \emph{producto exterior}; todas sus filas son múltiplos de $s^{\top}$. La forma (iii) descompone un producto de matrices en suma de productos exteriores y es la base de la descomposición en valores singulares y de las aproximaciones de bajo rango del módulo~3. En el \cref{ej:intercambio}, $a_1s_1^{\top}+a_2s_2^{\top}=\begin{pmatrix}0&1\\0&3\end{pmatrix}+\begin{pmatrix}2&0\\4&0\end{pmatrix}=AB$.

\section{Propiedades algebraicas}

\begin{proposicion}\label{prop:algebra}
Siempre que los tamaños permitan formar los productos:
\begin{enumerate}
  \item (asociativa) $(AB)C=A(BC)$;
  \item (distributivas) $A(B+C)=AB+AC$ y $(B+C)A=BA+CA$;
  \item (identidad) $I_mA=A=AI_n$ si $A\in\R^{m\times n}$;
  \item (traspuesta) $(AB)^{\top}=B^{\top}A^{\top}$.
\end{enumerate}
\end{proposicion}

La distributiva por la izquierda y la distributiva por la derecha son propiedades distintas: no puede pasarse de una a otra cambiando el lado (ejercicio~\ref{eje:4-distributiva}), porque el producto no es conmutativo.

\begin{observacion}[No conmutatividad]\label{obs:no-conmuta}
En general $AB\neq BA$, incluso si ambos productos están definidos y son del mismo tamaño. Sean $R=\begin{pmatrix}0&-1\\1&0\end{pmatrix}$ (rotación de $90^\circ$), $S=\begin{pmatrix}2&0\\0&0{,}5\end{pmatrix}$ (escalado) y $x=(3,1)$. Escalar y después rotar da $R(Sx)=R(6;\,0{,}5)=(-0{,}5;\,6)$; rotar y después escalar da $S(Rx)=S(-1,3)=(-2;\,1{,}5)$. El orden de las transformaciones importa.
\end{observacion}

\begin{definicion}[Matriz ortogonal]\label{def:ortogonal}
Una matriz $Q\in\R^{n\times n}$ es \emph{ortogonal} si $Q^{\top}Q=I_n$, es decir, si sus columnas son unitarias y ortogonales dos a dos.
\end{definicion}

\begin{proposicion}\label{prop:ortogonal}
Si $Q$ es ortogonal, entonces $(Qu)\cdot(Qv)=u\cdot v$ para todos $u,v\in\R^{n}$. En consecuencia, $Q$ conserva normas, distancias y ángulos.
\end{proposicion}

Las rotaciones $R_\theta$ son ortogonales (ejercicio~\ref{eje:4-ortogonal}), lo que generaliza la \cref{prop:rotacion}. Las codificaciones posicionales rotatorias (RoPE) del módulo~14 aplican precisamente rotaciones de este tipo a las consultas y las claves.

\section{Coste y matriz de Gram}

\begin{proposicion}[Coste del producto]\label{prop:coste}
Calcular $AB$ con $A\in\R^{m\times n}$ y $B\in\R^{n\times p}$ mediante la \cref{def:producto} requiere $mnp$ multiplicaciones y $mp(n-1)$ sumas.
\end{proposicion}

Aunque $(AB)C=A(BC)$, el coste de los dos órdenes puede ser muy distinto (ejercicio~\ref{eje:4-coste}). Elegir bien el orden de los paréntesis es una optimización habitual en el cálculo con tensores.

\begin{definicion}[Matriz de Gram]\label{def:gram}
Sea $X\in\R^{n\times d}$ la matriz cuyas filas son los vectores $x_1^{\top},\dots,x_n^{\top}$. Su \emph{matriz de Gram} es $G=XX^{\top}\in\R^{n\times n}$, con $G_{ij}=x_i\cdot x_j$.
\end{definicion}

La matriz de Gram es simétrica, porque $G^{\top}=(XX^{\top})^{\top}=XX^{\top}$, y su diagonal contiene las normas al cuadrado, $G_{ii}=\norm{x_i}^{2}$. Un único producto de matrices calcula los $n^{2}$ productos escalares entre $n$ vectores.

\addtocounter{definicion}{1}%

\begin{observacion}
En la atención, las consultas y las claves de los $n$ tokens de una secuencia se apilan como filas de matrices $Q,K\in\R^{n\times d_k}$, y todas las puntuaciones se calculan con un único producto $QK^{\top}\in\R^{n\times n}$, que generaliza la matriz de Gram. Las GPU están diseñadas para ejecutar este tipo de productos con gran rapidez (módulo~17).
\end{observacion}


% =====================================================================
\chapter{Embeddings, analogías y atención}
\label{tema:5}
% =====================================================================

Este tema aplica las herramientas anteriores a la entrada y la salida de un modelo de lenguaje. La capa de embeddings es una matriz; leer el vector de un token es multiplicar por un vector \emph{one-hot}; las relaciones semánticas aparecen, de forma aproximada, como direcciones del espacio; y la atención compara tokens mediante productos escalares. La presentación sigue el capítulo~5 de Jurafsky y Martin~\cite{jm3} y la clase~2 de CS224N~\cite{cs224n}, junto con los artículos de Mikolov y coautores~\cite{mikolov-a,mikolov-b} y de Vaswani y coautores~\cite{vaswani}.

\section{La matriz de embeddings}

En este tema se usa la notación del módulo~1: $\V$ es el vocabulario, $K=\abs{\V}$ su tamaño y los tokens se identifican con sus índices $0,1,\dots,K-1$. Las filas de las matrices indexadas por tokens se numeran también desde~0.

\begin{definicion}[Matriz de embeddings]\label{def:embedding}
Sea $d\ge 1$. Una \emph{matriz de embeddings} es una matriz $E\in\R^{K\times d}$. El \emph{embedding} del token $i$ es la fila $i$ de $E$, que se denota $E_{i,:}\in\R^{d}$ (escrita como columna cuando se opera con ella).
\end{definicion}

\begin{definicion}[Vector one-hot]
El vector \emph{one-hot} del token $i$ es el vector $e_i\in\R^{K}$ que tiene un~1 en la posición $i$ y ceros en el resto.
\end{definicion}

\begin{proposicion}\label{prop:onehot}
Para todo token $i$, $e_i^{\top}E=E_{i,:}^{\top}$; es decir, multiplicar el vector one-hot de un token por la matriz de embeddings selecciona su fila. Si $x_{1:n}$ es una secuencia de tokens y $O\in\R^{n\times K}$ es la matriz cuyas filas son $e_{x_1}^{\top},\dots,e_{x_n}^{\top}$, entonces $X=OE\in\R^{n\times d}$ es la matriz cuya fila $t$ es el embedding de $x_t$.
\end{proposicion}

\addtocounter{definicion}{1}%

\begin{observacion}[Indexar en lugar de multiplicar]
El producto $e_i^{\top}E$ requiere $Kd$ multiplicaciones, casi todas por cero; leer directamente la fila $i$ solo copia $d$ números. Las bibliotecas de redes neuronales implementan la capa de embeddings como una indexación, aunque matemáticamente sea el producto de la \cref{prop:onehot}. En GPT-2, con $K=50\,257$~\cite[§2.3]{radford} y $d=768$, la matriz $E$ tiene $50\,257\cdot 768=38\,597\,376$ parámetros, unos $154$~MB en precisión de 32 bits.
\end{observacion}

\section{De dónde salen los embeddings}

Los embeddings no se diseñan a mano como en el \cref{ej:juguete}: se inicializan al azar y se ajustan por descenso de gradiente (módulos~4 y~7) para que el modelo prediga bien el token siguiente. La justificación de que este procedimiento produzca vectores con significado es la \emph{hipótesis distribucional}: las palabras que aparecen en contextos parecidos tienden a tener significados parecidos~\cite{harris}. Firth la resumió en la frase «you shall know a word by the company it keeps»~\cite{firth}. Dos tokens que aparecen en contextos similares reciben gradientes similares y acaban con embeddings próximos en el sentido del coseno~\cite[cap.~5]{jm3}.

Antes de los Transformers, los embeddings se entrenaban por separado con modelos como word2vec~\cite{mikolov-a}, que aprende un vector por palabra prediciendo las palabras vecinas en una ventana. En un LLM, en cambio, la matriz $E$ se aprende junto con el resto de los parámetros.

\section{La capa de salida y el atado de pesos}

Al final del modelo, el vector $h\in\R^{d}$ de la última posición se compara con un vector por cada token del vocabulario. Si $U\in\R^{K\times d}$ es la matriz de salida, los logits son
\[
z=Uh\in\R^{K},\qquad z_i=U_{i,:}\cdot h,
\]
y la softmax del módulo~1 (tema~4) los convierte en una distribución sobre el token siguiente. Cada logit es, por tanto, un producto escalar entre $h$ y un vector asociado al token.

\begin{definicion}[Atado de pesos]
Se dice que un modelo \emph{ata} los pesos de entrada y salida (\emph{weight tying}) si usa $U=E$: la misma matriz sirve para leer los embeddings de entrada y para calcular los logits.
\end{definicion}

Press y Wolf~\cite{press} mostraron que esta elección mejora los modelos de lenguaje recurrentes, y el Transformer original la adopta~\cite[§3.4]{vaswani}. Ahorra $Kd$ parámetros, que en GPT-2 son $38\,597\,376$, y da a ambos usos una interpretación común: el logit del token $i$ es alto cuando $h$ apunta en la dirección del embedding de $i$.

\section{Analogías y aritmética de vectores}

Sea $v_w$ el embedding de la palabra $w$. Si un rasgo como el género corresponde aproximadamente a una dirección del espacio, la diferencia $v_{\text{mujer}}-v_{\text{hombre}}$ debería parecerse a $v_{\text{reina}}-v_{\text{rey}}$. Mikolov, Yih y Zweig~\cite{mikolov-b} observaron este fenómeno en embeddings aprendidos y propusieron resolver analogías con el método siguiente.

\begin{definicion}[Analogía por desplazamiento de vectores]\label{def:analogia}
Dada la pregunta «$a$ es a $b$ como $?$ es a $c$», se calcula $q=v_a-v_b+v_c$ y se responde con
\[
\hat w=\operatorname*{arg\,max}_{w\in\V\setminus\{a,b,c\}}\cos(v_w,q).
\]
\end{definicion}

\begin{proposicion}\label{prop:paralelogramo}
Si $v_a-v_b=v_{w^{*}}-v_c$, entonces $q=v_{w^{*}}$ y $\cos(v_{w^{*}},q)=1$, de modo que $w^{*}$ es una respuesta de máxima similitud.
\end{proposicion}

\addtocounter{definicion}{1}%

\begin{observacion}[Limitaciones]
Con embeddings reales, la igualdad de la \cref{prop:paralelogramo} solo se cumple de forma aproximada, y $q$ suele quedar más cerca de $v_a$ que de la respuesta esperada; la exclusión de $a$, $b$ y $c$ en la \cref{def:analogia} es esencial para que el método «funcione». Nissim, van Noord y van der Goot~\cite{nissim} analizan cómo esta exclusión distorsiona las conclusiones que se extraen de las analogías. En el cuaderno de prácticas del módulo, con 40 palabras en 16 dimensiones y ruido, la consulta rey $-$ hombre $+$ mujer da «reina» con coseno $0{,}924$, pero sin la exclusión «rey» queda segunda con $0{,}835$. Además, los embeddings heredan los sesgos de los textos de entrenamiento: Bolukbasi y coautores~\cite{bolukbasi} encontraron analogías que reproducen estereotipos de género. La idea de que los conceptos son direcciones reaparece en la interpretabilidad de los LLM (módulo~23).
\end{observacion}

\section{Adelanto: el producto escalar en la atención}

La atención (módulo~13) es un mecanismo de búsqueda por similitud. Cada token emite una consulta $q\in\R^{d_k}$, la compara por producto escalar con las claves $k_1,\dots,k_n$ de los tokens anteriores y reparte su atención con una softmax:
\[
s_j=\frac{q\cdot k_j}{\sqrt{d_k}},\qquad \alpha=\softmax(s_1,\dots,s_n).
\]
Para todos los tokens a la vez, las puntuaciones forman la matriz $QK^{\top}/\sqrt{d_k}$ (\cref{def:gram}).

\addtocounter{definicion}{1}%

El factor $1/\sqrt{d_k}$ tiene una justificación estadística~\cite[§3.2.1]{vaswani}: si las componentes de $q$ y $k$ son variables independientes de media~0 y varianza~1, entonces $\Var(q\cdot k)=d_k$ (ejercicio~\ref{eje:5-atencion}). Sin el factor, para $d_k$ grande las puntuaciones tendrían valores muy dispares, la softmax quedaría casi concentrada en un solo token y los gradientes serían casi nulos. Dividir por $\sqrt{d_k}$ devuelve la varianza a~1.



\backmatter
% =====================================================================
\begin{thebibliography}{99}
\addcontentsline{toc}{chapter}{Bibliografía}
\markboth{Bibliografía}{}
\small

\bibitem{bhk}
A.~Blum, J.~Hopcroft y R.~Kannan.
\newblock \emph{Foundations of Data Science}.
\newblock Cambridge University Press, 2020.
\newblock Cap.~2, «High-Dimensional Space». Texto del curso CS~4850 de Cornell.
\newblock \url{https://www.cs.cornell.edu/jeh/book.pdf}

\bibitem{bolukbasi}
T.~Bolukbasi, K.-W. Chang, J.~Zou, V.~Saligrama y A.~Kalai.
\newblock Man is to computer programmer as woman is to homemaker? Debiasing word embeddings.
\newblock En \emph{Advances in Neural Information Processing Systems 29}, 2016.
\newblock \url{https://arxiv.org/abs/1607.06520}

\bibitem{vmls}
S.~Boyd y L.~Vandenberghe.
\newblock \emph{Introduction to Applied Linear Algebra: Vectors, Matrices, and Least Squares}.
\newblock Cambridge University Press, 2018.
\newblock Texto del curso EE103 de Stanford.
\newblock \url{https://web.stanford.edu/~boyd/vmls/}

\bibitem{brown}
T.~B. Brown \emph{et al.}
\newblock Language models are few-shot learners.
\newblock En \emph{Advances in Neural Information Processing Systems 33}, 2020.
\newblock \url{https://arxiv.org/abs/2005.14165}

\bibitem{mml}
M.~P. Deisenroth, A.~A. Faisal y C.~S. Ong.
\newblock \emph{Mathematics for Machine Learning}.
\newblock Cambridge University Press, 2020.
\newblock Cap.~2, «Linear Algebra», y cap.~3, «Analytic Geometry».
\newblock \url{https://mml-book.github.io/}

\bibitem{firth}
J.~R. Firth.
\newblock A synopsis of linguistic theory, 1930--1955.
\newblock En \emph{Studies in Linguistic Analysis}, págs.~1--32. Blackwell, Oxford, 1957.

\bibitem{harris}
Z.~S. Harris.
\newblock Distributional structure.
\newblock \emph{Word}, 10(2--3):146--162, 1954.

\bibitem{jm3}
D.~Jurafsky y J.~H. Martin.
\newblock \emph{Speech and Language Processing}, 3.ª ed., borrador del 19 de agosto de 2026.
\newblock Cap.~5, «Embeddings».
\newblock \url{https://web.stanford.edu/~jurafsky/slp3/}

\bibitem{mikolov-a}
T.~Mikolov, K.~Chen, G.~Corrado y J.~Dean.
\newblock Efficient estimation of word representations in vector space.
\newblock arXiv:1301.3781, 2013.
\newblock \url{https://arxiv.org/abs/1301.3781}

\bibitem{mikolov-b}
T.~Mikolov, W.-t. Yih y G.~Zweig.
\newblock Linguistic regularities in continuous space word representations.
\newblock En \emph{Proceedings of NAACL-HLT}, págs.~746--751, 2013.
\newblock \url{https://aclanthology.org/N13-1090/}

\bibitem{nissim}
M.~Nissim, R.~van Noord y R.~van der Goot.
\newblock Fair is better than sensational: man is to doctor as woman is to doctor.
\newblock \emph{Computational Linguistics}, 46(2):487--497, 2020.

\bibitem{press}
O.~Press y L.~Wolf.
\newblock Using the output embedding to improve language models.
\newblock En \emph{Proceedings of EACL}, vol.~2, págs.~157--163, 2017.
\newblock \url{https://aclanthology.org/E17-2025/}

\bibitem{radford}
A.~Radford, J.~Wu, R.~Child, D.~Luan, D.~Amodei e I.~Sutskever.
\newblock Language models are unsupervised multitask learners.
\newblock Informe técnico, OpenAI, 2019.

\bibitem{strang}
G.~Strang.
\newblock \emph{Introduction to Linear Algebra}, 4.ª ed.
\newblock Wellesley-Cambridge Press, 2009.
\newblock Texto de los cursos 18.06 y 18.06SC del MIT. Hay 5.ª ed. (2016) y 6.ª ed. (2023); la numeración de secciones y problemas citada es la de la 4.ª ed.

\bibitem{mit1806}
G.~Strang.
\newblock 18.06: Linear Algebra.
\newblock Massachusetts Institute of Technology, primavera de 2010. Hoja de problemas~1.
\newblock \url{https://web.mit.edu/18.06/www/Spring10/}

\bibitem{mit1806sc}
G.~Strang.
\newblock 18.06SC: Linear Algebra.
\newblock Massachusetts Institute of Technology, MIT OpenCourseWare, otoño de 2011.
\newblock \url{https://ocw.mit.edu/courses/18-06sc-linear-algebra-fall-2011/}

\bibitem{vaswani}
A.~Vaswani, N.~Shazeer, N.~Parmar, J.~Uszkoreit, L.~Jones, A.~N. Gomez, Ł.~Kaiser e I.~Polosukhin.
\newblock Attention is all you need.
\newblock En \emph{Advances in Neural Information Processing Systems 30}, 2017.
\newblock \url{https://arxiv.org/abs/1706.03762}

\bibitem{cs224n}
D.~Yang y Y.~Choi.
\newblock CS224N: Natural Language Processing with Deep Learning.
\newblock Stanford University, invierno de 2026. Clase~2, «Word Vectors».
\newblock \url{https://web.stanford.edu/class/cs224n/}

\bibitem{d2l}
A.~Zhang, Z.~C. Lipton, M.~Li y A.~J. Smola.
\newblock \emph{Dive into Deep Learning}, versión 1.0.3.
\newblock Cambridge University Press, 2023.
\newblock \url{https://d2l.ai/}

\end{thebibliography}


\end{document}
